\section{The on-shell scattering operator as a direct integral}
\label{sec:scattering-direct-integral}
% BEGIN CURATED SUBJECT INDEX
\index{direct integral}
\index{wave operator}
\index{scattering matrix}
\index{scattering!on-shell}
% END CURATED SUBJECT INDEX

The energy representation constructed in
Section~\ref{subsec:free-direct-integral} becomes a scattering representation
only after the M{\o}ller operators have been defined.  We now prove the missing
step.  This is also the point at which the symbol $\Scattering$ acquires its
full operator meaning; it is not an unspecified object that happens to
commute with a Hamiltonian.

\subsection{Intertwining is stronger than a formal commutator}
% BEGIN CURATED SUBJECT INDEX
\index{wave operator}
\index{scattering matrix}
\index{boundary condition!incoming and outgoing}
% END CURATED SUBJECT INDEX

Use the convention of Eq.~\eqref{eq:moller-wave-operators}:
$\Omega^{(+)}$ describes an incoming state and $\Omega^{(-)}$ an outgoing
state.  Suppose the strong limits exist on
$\Hil_{0,\mathrm{ac}}=\vsub{P}{ac}(H_0)\Hil$.  Translation of the limiting
time gives, first for the unitary groups and then for every bounded Borel
function $f$, the intertwining relation
\begin{equation}
  f(H)\Omega^{(\pm)}
  =\Omega^{(\pm)}f(H_0).
  \label{eq:wave-operator-borel-intertwining}
\end{equation}
The statement includes spectral projections, not only vectors in the domain
of $H_0$.  In particular,
\begin{equation}
  P_H(\Delta)\Omega^{(\pm)}
  =\Omega^{(\pm)}P_{H_0}(\Delta)
  \label{eq:wave-operator-spectral-intertwining}
\end{equation}
for every Borel set $\Delta$.

Define
\begin{equation}
  \Scattering=\Omega^{(-)\dagger}\Omega^{(+)}
  \quad\text{on }\Hil_{0,\mathrm{ac}}.
  \label{eq:scattering-operator-direct-integral}
\end{equation}
Using Eq.~\eqref{eq:wave-operator-spectral-intertwining} twice gives
\begin{equation}
  P_{H_0}(\Delta)\Scattering
  =\Scattering P_{H_0}(\Delta).
  \label{eq:S-strongly-commutes-H0}
\end{equation}
This is \term{strong commutation} with $H_0$.  The shorthand
$[\Scattering,H_0]=0$ is safe only after
Eq.~\eqref{eq:S-strongly-commutes-H0} has been established; a formal
commutator on an unspecified common domain would be weaker.

If the wave operators are isometries, $\Scattering$ is an isometry.  If their
ranges coincide with the absolutely continuous subspace of $H$---the
asymptotic-completeness statement---then $\Scattering$ is unitary on
$\Hil_{0,\mathrm{ac}}$.  Completeness is therefore the step that upgrades
flux conservation in all represented channels to operator unitarity.

\subsection{Decomposability and the on-shell fibers}
% BEGIN CURATED SUBJECT INDEX
\index{measure theory}
\index{spectral measure}
\index{decomposable operator}
\index{commutant}
\index{scattering matrix}
\index{scattering!on-shell}
% END CURATED SUBJECT INDEX

Let $U_0$ be a multiplicity representation of $H_0$,
\begin{equation}
  U_0\Hil_{0,\mathrm{ac}}
  =\int_{\acspectrum(H_0)}^\oplus
    \mathcal K_E\dd\mu_0(E),
  \qquad
  (U_0H_0U_0^\dagger\psi)(E)=E\psi(E).
  \label{eq:H0-multiplicity-stage-II}
\end{equation}
Here $\acspectrum(H_0)$ denotes an essential support of the
absolutely continuous spectral measure; changing it on a null set changes no
direct-integral vector.
Equation~\eqref{eq:S-strongly-commutes-H0} says exactly that
$U_0\Scattering U_0^\dagger$ lies in the commutant of the diagonal algebra.
The decomposability theorem of Section~\ref{subsec:direct-integrals} therefore
yields
\begin{equation}
  U_0\Scattering U_0^\dagger
  =\int_{\acspectrum(H_0)}^\oplus
    \Scattering(E)\dd\mu_0(E).
  \label{eq:S-direct-integral}
\end{equation}
The family $E\mapsto\Scattering(E)$ is measurable and essentially bounded;
when $\Scattering$ is unitary, $\Scattering(E)$ is unitary for
$\mu_0$-almost every $E$.  This almost-everywhere family is the rigorous
meaning of the \term{on-shell scattering matrix}.

The conclusion is representation independent, although its displayed matrix
is not.  Replacing the measure by
$\rmd\widetilde\mu(E)=w(E)\rmd\mu_0(E)$ rescales fiber coordinates by
$w(E)^{-1/2}$.  Channel amplitudes and delta functions acquire compensating
density factors.  Probabilities, pole positions after analytic continuation,
and the unitary equivalence class of the decomposable operator do not change.

\subsection{The free particle and partial waves}
% BEGIN CURATED SUBJECT INDEX
\index{direct integral!fiber}
\index{decomposable operator}
\index{partial-wave expansion}
\index{scattering matrix}
% END CURATED SUBJECT INDEX

For the three-dimensional free particle, Eq.~\eqref{eq:free-energy-unitary}
gives
\begin{equation}
  \mathcal K_E=L^2(S^2,\rmd\Omega),
  \qquad E>0.
  \label{eq:free-fiber-angular}
\end{equation}
Thus $\Scattering(E)$ is an operator on angular wave functions at fixed
energy.  Rotational invariance implies that it commutes with the angular
momentum representation.  For a spinless central interaction,
\begin{equation}
  \Scattering(E)Y_{lm}
  =S_l(E)Y_{lm},
  \qquad
  S_l(E)=\rme^{2\rmi\delta_l(E)}.
  \label{eq:S-fiber-partial-waves}
\end{equation}
The familiar sequence of partial-wave numbers is therefore a second,
discrete decomposition inside one energy fiber.  Energy is the direct-
integral base; $(l,m)$ are multiplicity coordinates in the fiber.

In generalized momentum notation the same fact is often written
\begin{equation}
  \bra{\bm k'}\Scattering\ket{\bm k}
  =\delta(E_{k'}-E_k)\,
   \mathsf s_E(\widehat{\bm k}',\widehat{\bm k}) .
  \label{eq:S-kernel-energy-delta}
\end{equation}
The energy delta function is not an additional dynamical assumption.  It is
the distributional kernel of a decomposable operator.  Its normalization
depends on the chosen momentum measure, which is why working first with
packets avoids spurious factors.

\subsection{Thresholds and channel multiplicity}
% BEGIN CURATED SUBJECT INDEX
\index{measure theory}
\index{resolvent}
\index{direct integral}
\index{decomposable operator}
\index{Jost function}
\index{wave operator}
\index{scattering matrix}
\index{scattering!on-shell}
\index{scattering!multichannel}
\index{exact WKB}
% END CURATED SUBJECT INDEX

Suppose channel $a$ has threshold $E_a$.  A convenient fiber is
\begin{equation}
  \mathcal K_E
  =\bigoplus_{a:E>E_a}\mathcal K_{a,E}.
  \label{eq:open-channel-fiber}
\end{equation}
The fiber dimension jumps at thresholds.  Since isolated threshold energies
form a null set for Lebesgue measure, the measurable direct integral has no
difficulty with the jump.  Analytic continuation does: the channel momentum
\begin{equation}
  k_a(E)=\frac{\sqrt{2\mu_a(E-E_a)}}{\hbar}
  \label{eq:channel-momentum-threshold}
\end{equation}
introduces a branch point and a sheet choice.  Direct-integral theory tells us
what the real-axis on-shell operator is almost everywhere; it does not supply
the analytic relation between different energies or sheets.  Jost functions,
resolvent boundary values, and exact WKB provide that extra structure.

\begin{checkpointbox}
The statement $[\Scattering,H_0]=0$ has now been unpacked into three claims:
wave operators exist, their spectral projections intertwine, and the
resulting scattering operator is decomposable in the spectral representation
of the \textit{comparison} Hamiltonian.  None of these claims says that
$\Scattering$ commutes with an arbitrary interacting Hamiltonian.
\end{checkpointbox}

\subsection{What a finite channel calculation approximates}
% BEGIN CURATED SUBJECT INDEX
\index{scattering matrix}
\index{CDCC}
\index{nuclear breakup}
% END CURATED SUBJECT INDEX

A finite coupled-channel matrix $S^{(N)}(E)$ approximates the fiber operator
$\Scattering(E)$, not the full operator before energy decomposition.  Its
truncation contract must specify the retained channel subspace
$P_N(E)\mathcal K_E$, the energy interval, and the norm or observable in which
convergence is tested.  Pointwise convergence away from thresholds need not
be uniform at a threshold or pole.  Conversely, exact unitarity of a finite
matrix does not prove that omitted breakup channels are negligible.  The CDCC
Stage will turn these observations into a numerical convergence protocol.
